\section*{Erd\H{o}s Problem 1026}

\paragraph{Statement and precise scope.}
For a sequence of distinct real numbers, consider the largest sum of a nonempty increasing or decreasing subsequence. The original request to ``determine'' this quantity has more than one possible interpretation. We give an exact recurrence for any given sequence and prove the sharp universal asymptotic constant in the clarified question:
\[
 M(x_1,\ldots,x_n)\geq\frac1{\sqrt n}\sum_{i=1}^n x_i.
 \tag{1}
\]
For positive sequences of square length \(n=k^2\) with total sum one, this proves the requested lower bound \(1/k\); no constant larger than one can replace the coefficient in (1) uniformly as \(n\to\infty\). The sharper extremal constants for individual nonsquare lengths are separate from the assertion proved here.

\paragraph{The exact value for any input sequence.}
Let \(I_i\) and \(D_i\) be the largest sums of increasing and decreasing subsequences ending at position \(i\), respectively. Then
\[
 I_i=x_i+\max\bigl(\{0\}\cup\{I_j:j<i,\ x_j<x_i\}\bigr),
\]
\[
 D_i=x_i+\max\bigl(\{0\}\cup\{D_j:j<i,\ x_j>x_i\}\bigr). \tag{2}
\]
Every such subsequence either consists of its final entry or extends one ending at an eligible earlier entry. Conversely each choice in (2) produces such a subsequence. Induction on \(i\) therefore proves the recurrence, including for negative entries. The desired value is exactly
\(M=\max_i\max(I_i,D_i)\). Computing all the maxima by scanning earlier positions takes \(O(n^2)\) comparisons and additions; retaining a maximizing predecessor recovers an actual subsequence.

\paragraph{A stronger weighted inequality for positive entries.}
Assume first that all \(x_i>0\), and write
\(I=\max I_i\), \(D=\max D_i\). Associate to position \(i\) the square
\[
 Q_i=(I_i-x_i,I_i]\times(D_i-x_i,D_i].
\]
It has area \(x_i^2\), and lies in \((0,I]\times(0,D]\), since both recurrences include the option zero.

These squares are pairwise disjoint. Indeed, for \(i<j\), if \(x_i<x_j\), then (2) gives \(I_j-x_j\geq I_i\), so the two horizontal intervals are disjoint. If \(x_i>x_j\), then \(D_j-x_j\geq D_i\), so the two vertical intervals are disjoint. The half-open convention handles equality at a boundary without double counting.

Adding areas now gives
\[
 ID\geq\sum_i x_i^2,
 \qquad M^2\geq ID\geq\sum_i x_i^2
 \geq\frac1n\left(\sum_i x_i\right)^2. \tag{3}
\]
The last inequality is Cauchy--Schwarz. Taking nonnegative square roots proves (1). In fact the intermediate bound \(M\geq(\sum x_i^2)^{1/2}\) is stronger than the normalized-sum assertion.

\paragraph{Allowing arbitrary real entries.}
If some entries are positive, restrict to the subsequence of positive entries. It has length \(m\leq n\) and positive sum \(S_+\). If the full sum \(S\) is positive, then (3) gives
\[
 M\geq S_+/\sqrt m\geq S/\sqrt n.
\]
If \(S\leq0\), any positive singleton already proves (1). Finally, if every entry is nonpositive, the largest entry is at least the average \(S/n\), and
\(S/n\geq S/\sqrt n\). Taking that singleton proves (1) in the remaining case. Thus no positivity assumption is silently imposed on the original sequence.

\paragraph{Sharpness of the asymptotic constant.}
Let \(n=k^2\), and consider the permutation formed from \(k\) consecutive decreasing blocks,
\[
 (k,k-1,\ldots,1),\ (2k,2k-1,\ldots,k+1),\ \ldots,
 (k^2,k^2-1,\ldots,k^2-k+1).
\]
Call its entries \(r_i\). An increasing subsequence takes at most one entry from each block, and a decreasing subsequence lies in one block. Hence every monotone subsequence has at most \(k\) entries.

For any \(\eta>0\), set \(x_i=1+\eta r_i/k^2\). These are distinct positive reals, with exactly the displayed order relations and all lying in \((1,1+\eta]\). Therefore
\[
 M\leq k(1+\eta),\qquad \sum_i x_i\geq k^2,
 \qquad \frac{\sqrt n\,M}{\sum_i x_i}\leq1+\eta.
\]
Let \(k\to\infty\) and \(\eta\to0\). This rules out every universal asymptotic coefficient greater than one. Normalizing each example by its total sum also shows that the square-length bound \(1/k\) is the infimum, even with distinct positive entries.

\paragraph{Attribution and assessment.}
The weighted conclusion appears in work of Tidor, Wang, and Yang and is implicit in work of Wagner. The square-packing exposition above follows the argument attributed to Aristotle in the discussion, described there by llllvvuu; KoishiChan supplies an alternative blow-up proof, and Cambie supplies the near-equal block examples. Sources: \url{https://www.erdosproblems.com/1026} and \url{https://www.erdosproblems.com/forum/thread/1026}, checked 24 September 2026.

The exact input recurrence, universal bound, square-length conjecture, and asymptotic constant are completely proved here. No external solution theorem or local Lean verification is assumed, and no new result is claimed.

\textbf{Completion rate estimate: 100\%.}
